\section{Theory}
\label{sec:theory}

This section states the physical model that \alibz{} uses to turn a laser-induced
breakdown spectrum into an estimate of elemental composition, plasma temperature
$T$, and electron density $\nel$. The development is deliberately explicit about
one point of discipline that governs the whole code base: the forward model that
predicts a spectrum and the inverse model that fits one share a single line-intensity
convention, implemented once and called from both directions \citep{alibz}. Where a
textbook relation is simplified, replaced by a data-driven anchor, or deliberately
omitted, we say so here and revisit it as a named assumption in
Section~\ref{sec:assumptions}.

Throughout, $z$ indexes the ionization stage of an element (\mbox{$z=1$} neutral,
$z=2$ singly ionized, and so on), $k$ an atomic energy level, and $j$ an individual
emission line. We write $\lambda_j$ for the air wavelength of line $j$, $g_k$ for the
statistical weight of level $k$, $E_k$ for its energy above the ground state of its
stage, and $A_{ki}$ for the Einstein spontaneous-emission coefficient of the
$k\!\rightarrow\!i$ transition. Atomic databases tabulate the product
$g_kA_{ki}$, which we abbreviate $gA_j$, so that quantity rather than $A_{ki}$ alone
appears below. Energies are in electron-volts, wavelengths in nanometers, and $T$
in kelvin unless stated otherwise.

\subsection{Line emission in local thermodynamic equilibrium}
\label{sec:theory-lte}

We assume (Assumption~A1) that, within a single ionization stage, the level
populations follow a Boltzmann distribution at one excitation temperature $T$. The
number density of emitters in upper level $k$ of stage $z$ is then
%
\begin{equation}
  \label{eq:boltzmann}
  N_k = N_z\,\frac{g_k}{U_z(T)}\,\exp\!\left(-\frac{E_k}{\kB T}\right),
\end{equation}
%
where $N_z$ is the total number density of the element in stage $z$, $\kB$ is the
Boltzmann constant, and $U_z(T)$ is the internal partition function of the stage,
%
\begin{equation}
  \label{eq:partition}
  U_z(T) = \sum_k g_k\,\exp\!\left(-\frac{E_k}{\kB T}\right).
\end{equation}
%
\alibz{} does not read tabulated partition functions. It evaluates
Equation~\eqref{eq:partition} directly over the union of lower and upper levels that
appear in the line list it holds for stage $z$ (module \texttt{sahaboltzmann}). This
couples the partition sum to the completeness of the line list and is a known
limitation (Assumption~A5); critically evaluated partition tables
\citep{irwin1981partition,barklem2016partition} are a planned replacement, and the salvaged
\texttt{dev1} prototype's spline-interpolated $(T,U)$ tables with Griem truncation
are recorded but not adopted.

For an optically thin, homogeneous source, the energy radiated per unit time and
solid angle in line $j$ is $I_j = N_k A_{ki}\,hc/\lambda_j$, with $h$ Planck's
constant and $c$ the speed of light \citep{griem1997principles,kunze2009plasma}; the
atomic constants are taken from the NIST Atomic Spectra Database \citep{kramida2023asd}.
Substituting Equation~\eqref{eq:boltzmann} and folding in the fraction $f_z(T,\nel)$
of the element in stage $z$ (Section~\ref{sec:theory-saha}) gives the quantity that
\alibz{} actually computes, the \emph{line emissivity} per unit element number
density:
%
\begin{equation}
  \label{eq:emissivity}
  \varepsilon_j(T,\nel) =
    \frac{1}{4\pi}\,\frac{hc}{\lambda_j}\,
    \frac{gA_j}{U_z(T)}\,
    \exp\!\left(-\frac{E_{k,j}}{\kB T}\right)\,
    f_z(T,\nel).
\end{equation}
%
Equation~\eqref{eq:emissivity} is the single intensity convention of the code
(\texttt{sahaboltzmann.line\_emissivity}). Two features are load-bearing. First, the
emissivity is in \emph{energy} units: the photon-energy factor $hc/\lambda_j$ is
present, so a blue line and a red line of equal $gA\exp(-E_k/\kB T)$ do not emit
equal measured energy. Second, the same function is called by the forward generator
and by the inverse solver, so the two cannot drift apart. We note this because an
earlier inverse model omitted the $hc/\lambda$ factor; because the Z300 spans
roughly \SIrange{190}{950}{\nano\meter}, that omission introduced a measured
blue-to-red bias of up to $\sim$\num{5}$\times$ in relative line weights, which the
shared convention removes.

The predicted area of line $j$ for an element at number density $N_\mathrm{el}$,
as it reaches the detector, is
%
\begin{equation}
  \label{eq:forward-line}
  Y_j = R(\lambda_j)\,N_\mathrm{el}\,\varepsilon_j(T,\nel),
\end{equation}
%
where $R(\lambda)$ is the instrument's spectral response (Assumption~A6). The
explicit-stage forward generator adds a positive free--free continuum and nuisance
emission from entrained air or argon; free--bound (recombination) continuum is not
modeled (Assumption~A5). Because Equation~\eqref{eq:emissivity} is in energy units,
$R(\lambda)$ must be an energy response; only its \SI{620}{\nano\meter} gain step is
currently measured and applied, and the ultraviolet-to-visible relative response is
uncalibrated (Assumption~A6).

\subsection{Ionization balance and electron density}
\label{sec:theory-saha}

The stage fractions $f_z$ in Equation~\eqref{eq:emissivity} follow from the Saha
equation \citep{saha1920ionization,griem1997principles} for successive stages in
equilibrium,
%
\begin{equation}
  \label{eq:saha}
  \frac{n_{z+1}}{n_z} =
    \frac{2}{\nel}\,\frac{U_{z+1}(T)}{U_z(T)}
    \left(\frac{2\pi m_\mathrm{e}\kB T}{h^2}\right)^{3/2}
    \exp\!\left(-\frac{\chi_z}{\kB T}\right),
\end{equation}
%
with $m_\mathrm{e}$ the electron mass, $\chi_z$ the ionization energy connecting
stage $z$ to $z+1$, and the factor $2$ the statistical weight of the free electron.
\alibz{} evaluates Equation~\eqref{eq:saha} through the thermal de~Broglie wavelength
$\Lambda = h/\sqrt{2\pi m_\mathrm{e}\kB T}$ (module \texttt{sahaboltzmann}), normalizes
across the stages present to obtain
%
\begin{equation}
  \label{eq:stage-fraction}
  f_z(T,\nel) = \frac{n_z}{\sum_{z'} n_{z'}},
\end{equation}
%
and passes $f_z$ into Equation~\eqref{eq:emissivity}. The electron density enters
only as $\log_{10}(\nel/\si{\per\cubic\centi\meter})$ internally.

Two physical corrections to Equation~\eqref{eq:saha} are standard in the dense-plasma
literature but are \emph{not} implemented, and we flag them here rather than imply
their presence. Ionization-potential depression (continuum lowering)
\citep{ecker1963lowering,stewart1966lowering} would reduce
each $\chi_z$ by a Debye term $\Delta\chi_z = (z+1)e^2/4\pi\varepsilon_0\lambda_\mathrm{D}$;
at LIBS conditions ($\nel\approx\SI{e17}{\per\cubic\centi\meter}$,
$T\approx\SI{8}{\kilo\kelvin}$, Debye length $\lambda_\mathrm{D}\approx\SI{20}{\nano\meter}$)
this is $\approx\SI{0.07}{\electronvolt}$ per unit charge, which changes the Saha
factors of the first two stages by roughly \SIrange{10}{25}{\percent}; the code
lowers no ionization energy. % src: Debye–Hückel estimate, reports/2026-11-03-lit-ingest-theory.md The McWhirter criterion
$\nel \gtrsim \SI{1.6e12}{}\,\sqrt{T}\,(\Delta E)^3$
\citep{mcwhirter1965spectral,griem1963validity,fujimoto1990validity}, a necessary lower
bound on
$\nel$ for the LTE assumption behind Equations~\eqref{eq:boltzmann}
and~\eqref{eq:saha} to hold, is likewise not evaluated. Both are recorded as
salvaged ideas awaiting A/B testing (Assumption~A1); satisfying the McWhirter bound is
in any case necessary but not sufficient for LTE in a transient LIBS plasma
\citep{cristoforetti2010lte}.

A decision specific to \alibz{} is that the stages are \emph{not} tied through
Equation~\eqref{eq:saha} in the inverse solve. The Saha fractions build the design
columns (Section~\ref{sec:theory-inverse}), but each (element, stage) density is an
independent non-negative unknown. There are two reasons. Analytically, $\nel$ is
unidentifiable from line amplitudes alone when the stages are freely tied, because a
per-stage density absorbs any common scale unless at least two stages of one element
carry design weight. Physically, the same element co-ablated from different mineral
phases ionizes differently, so forcing a single Saha ratio would misdescribe a
heterogeneous target. \alibz{} therefore reports the relative spread of an element's
independent per-stage estimates,
%
\begin{equation}
  \label{eq:stage-disagreement}
  D_\mathrm{el} = \frac{\max_z \hat N_z - \min_z \hat N_z}{\sum_z \hat N_z},
\end{equation}
%
as \texttt{stage\_disagreement}: $D_\mathrm{el}\!\approx\!0$ is consistent with a
single plasma, while $D_\mathrm{el}$ near $1$ flags a wrong plasma state, a non-LTE
condition, a blend-degenerate split, or phase heterogeneity (Assumption~A11).

Because the amplitude data alone constrain $\nel$ only weakly, \alibz{} measures it,
where possible, from line \emph{width}. Hydrogen Balmer lines are linear-Stark
dominated and carry an absolute density scale independent of any width calibration
\citep{gigosos1996hbeta,konjevic2002review}.
From the Stark full width at half maximum (FWHM) $\Delta\lambda_{\mathrm H}$ of
H$\alpha$, using the Gigosos--Gonz\'alez--Carde\~noso tables \citep{gigosos2003halpha},
\alibz{} sets
%
\begin{equation}
  \label{eq:halpha}
  \log_{10}\!\frac{\nel}{\si{\per\cubic\centi\meter}}
    = 17 + \frac{1}{\alpha}\,
      \log_{10}\!\frac{\Delta\lambda_{\mathrm H}}{\Delta\lambda_\mathrm{ref}},
\end{equation}
%
with $\alpha$ and $\Delta\lambda_\mathrm{ref}$ the fit constants (module
\texttt{stark}). On real samples this H$\alpha$ anchor returned
$\log_{10}\nel = 16.95$, $16.76$, and $17.32$ on three of four samples carrying a
usable H$\alpha$ line.

For non-hydrogenic lines the quadratic (impact) Stark width
\citep{griem1974broadening,griem1968semiempirical,konjevic1999nonhydrogenic} scales as
%
\begin{equation}
  \label{eq:stark-quad}
  w_j(\nel) \propto \frac{n_{\mathrm{eff},j}^4}{z_j^2}\,
    10^{\,\log_{10}\nel - 17},
  \qquad
  n_{\mathrm{eff},j}^2 = \frac{z_j^2\,\mathcal{R}_\infty}{\chi_z - E_{k,j}},
\end{equation}
%
where $n_{\mathrm{eff},j}$ is the effective principal quantum number of the upper
level and $\mathcal{R}_\infty$ the Rydberg energy. This channel is implemented
(module \texttt{stark}), with a ceiling $n_{\mathrm{eff}}\le 8$ to exclude
near-threshold Rydberg levels that would otherwise swamp the width objective, a
single global scale $c_4$ calibrated against a known density, and the instrumental
width intercept profiled analytically. It is nonetheless \emph{off by default},
because on this instrument class the predicted widths ($\sim$\SI{0.002}{\nano\meter}
at $\nel=\SI{e17}{\per\cubic\centi\meter}$) sit an order of magnitude below the
measured per-peak width scatter (\SIrange{0.03}{0.05}{\nano\meter}), so they are not
recoverable. In the absence of a usable H$\alpha$ line, \alibz{} adopts a default
prior $\log_{10}\nel \sim \mathcal{N}(17.0,\,0.5)$, searches inside the $1\sigma$ box,
and raises the flag \texttt{electron-density-at-bound} when the optimum rails against
it.

\subsection{Line profiles}
\label{sec:theory-profiles}

Each line is modeled as a Voigt profile, the convolution of a Gaussian of standard
deviation $\sigma$ (Doppler and instrumental broadening) with a Lorentzian of
half-width at half-maximum $\gamma$ (Stark and natural broadening). \alibz{} uses the
Olivero--Longbothum/Thompson approximation for the Voigt FWHM
\citep{olivero1977voigt,thompson1987voigt},
%
\begin{equation}
  \label{eq:voigt-fwhm}
  f_V = 0.5346\,f_L + \sqrt{0.2166\,f_L^2 + f_G^2},
  \qquad f_L = 2\gamma,\quad f_G = \sqrt{8\ln 2}\,\sigma,
\end{equation}
%
(module \texttt{voigt}), and evaluates the profile itself through the Faddeeva
function \citep{humlicek1982voigt,virtanen2020scipy}.
The thermal Doppler width is physically
$\sigma_\mathrm{D} = (\lambda/c)\sqrt{\kB T/m}$ for an emitter of mass $m$, but the
forward generator uses a \emph{fixed} width rather than tying $\sigma_\mathrm{D}$ to
$T$ and $m$ (Assumption~A9); grounding the Gaussian width in that relation would remove
one free parameter and is a recorded, untested option.

Two instrument-specific facts shape the profile model. The Z300 exports every
spectrum on a uniform \SI{0.0333}{\nano\meter} (\mbox{$1/30$}~nm) grid, which we have
shown to be a cubic-spline resample of coarser native pixels (native pitch
$\approx$ \SIrange{0.096}{0.198}{\nano\meter} across the three segments). Fitting on
the export grid biases fitted FWHM high by \SIrange{37}{80}{\percent}, biases line
area low by \SIrange{4}{6}{\percent}, and over-counts the information content roughly
four- to fivefold, because the spline introduces a lag-1 autocorrelation
$\rho(1)=0.92$ (falling as $0.76, 0.57, 0.40, 0.28, 0.20$ over lags $2$--$6$) that is
absent on the native grid, where the recovered noise is white (Assumption~A7).
\alibz{} therefore recovers and fits on the native grid, and integrates the model
over each native pixel rather than point-sampling it. The effective
point-spread function is handled per segment: an empirical PSF \citep{anderson2000epsf}
in the visible, a
spline-root location in the ultraviolet, and a point Gaussian in the near-infrared
(modules \texttt{peak\_location} and \texttt{epsf\_shapes}, with the per-segment
choice stored in a packaged shape file). A forward-convolved,
pixel-integrated line-spread function that would let the self-absorption parameters
of Section~\ref{sec:theory-sa} measure plasma physics rather than absorb instrument
shape is planned but not built; the interim mitigation is a robust \texttt{soft\_l1}
loss in the refinement fit.

\subsection{Self-absorption}
\label{sec:theory-sa}

Strong resonance lines are reabsorbed within the plasma, which flattens their cores
and understates the emitting population---exactly the population that quantification
needs. \alibz{} treats self-absorption with a homogeneous-slab escape factor
\citep{cowan1948selfabsorption,irons1979escape}. The
integrated optical depth of line $j$ scales with the lower-level population,
%
\begin{equation}
  \label{eq:kappa}
  \tau_j = s_\tau\,c_s\,\kappa_j,
  \qquad
  \kappa_j \propto \lambda_j^2\,gA_j\,
    \frac{\exp(-E_{i,j}/\kB T)}{U_z(T)}\,f_z(T,\nel),
\end{equation}
%
where $E_{i,j}$ is the \emph{lower}-level energy of the transition, $c_s$ the
species concentration, and $s_\tau$ a single global scale (module \texttt{absorption}).
The measured area is compressed by the escape factor
%
\begin{equation}
  \label{eq:escape}
  f_\mathrm{SA}(\tau) = \frac{1-e^{-\tau}}{\tau},
\end{equation}
%
which tends to $1$ as $\tau\!\rightarrow\!0$ (optically thin) and to $1/\tau$ for
$\tau\gg1$ (saturated). In the inverse solve the escape factor is applied by a
damped, Aitken-accelerated fixed point with a divergence guard: solve optically thin,
form $\tau_j$ from Equation~\eqref{eq:kappa} with $\kappa_j$ normalized to a frozen
reference line, compress the weights by Equation~\eqref{eq:escape}, and re-solve
(Section~\ref{sec:theory-inverse}). The pass-2 refinement caps $\tau\le8$; the
indexer's doublet channel caps $\tau\le50$.

Where two ground-state resonance lines of one species are resolved, their measured
area ratio inverts directly to an optical depth
\citep{elsherbini2005selfabsorption,sun2009internalreference,bulajic2002selfabsorption}. Both members share the lower level,
so their optical depths are locked in the ratio $s_r$ of their absorption strengths.
With $R_0$ the optically thin emissivity ratio of the strong to the weak member, the
observed area ratio is
%
\begin{equation}
  \label{eq:doublet}
  R(\tau) = R_0\,\frac{f_\mathrm{SA}(s_r\,\tau)}{f_\mathrm{SA}(\tau)},
\end{equation}
%
which \alibz{} solves for the weak-member $\tau$ (\texttt{invert\_doublet\_tau})
and uses to anchor that species' self-absorption (\texttt{\_freeze\_doublet\_taus},
gated off by default). The compression is real and large: the \mbox{K\,\textsc{i}}
\SI{766.5}{\nano\meter}/\SI{769.9}{\nano\meter} doublet has a thin ratio of $2.02$
but a measured ratio of $1.30$ on sample MW2-112, and the Na\,D doublet $2.00$ versus
a measured $1.51$---a $\sim$\SI{35}{\percent} amplitude compression that no thin
linear model can absorb.

When a resonance line is resolved well enough to fit its shape, \alibz{} uses a
cold-absorber model in which the emitted Voigt profile is attenuated by a
same-shaped, slightly shifted absorber,
%
\begin{equation}
  \label{eq:sa-voigt}
  I(x) = A\,V(x-\mu)\,
    \exp\!\left[-\tau_a\,\frac{V(x-\mu-\delta)}{V(0)}\right],
\end{equation}
%
with absorber depth $\tau_a$ and center shift $\delta$ (\texttt{refinement.sa\_voigt}).
This resolved-profile model is applied only when the transition has a resonance-capable
lower level ($E_i\le\SI{0.2}{\electronvolt}$). It separates two areas that must not be
confused: the \emph{observed area}, the integral of the attenuated profile actually
seen, which is stored in the peak table and is faithful to the data; and the
\emph{emission area} $A$, the thin-equivalent flux the line would emit with no
self-absorption, which is what Equation~\eqref{eq:forward-line} relates to
concentration and which is carried in the decision record. On MW2-112 the
emission-to-observed ratios ranged from $1.24$ to $7.0$ across fifteen merged lines.
The verdict is a non-destructive tag; downstream consumers must not re-apply a
self-absorption correction to a row that already carries \texttt{observed\_area}.

A stimulated-emission (induced-emission) correction to the net absorption,
%
\begin{equation}
  \label{eq:stim}
  1 - \exp\!\left(-\frac{hc}{\lambda\,\kB T}\right),
\end{equation}
%
is implemented (\texttt{stimulated\_emission\_factor}) and is wavelength-dependent
($\approx0.99$ in the ultraviolet versus $\approx0.89$ at the
\mbox{K\,\textsc{i}}~\SI{770}{\nano\meter} line at LIBS temperatures), so it cannot be
absorbed into the single global scale $s_\tau$. It is nonetheless \emph{off by
default}: an A/B test on 38 spectra found neutral fit quality (median
$\Delta r^2 = 0$) but composition swings exceeding \SI{20}{\percent} on $14$ of $38$
samples, because with only the doublet-anchored channel active the factor re-tilts a
near-degenerate alkali/self-absorption direction without adding constraint.

\subsection{Temperature from stage consistency}
\label{sec:theory-thermometer}

The amplitude objective is nearly flat in the plasma state, a fact that forces the
choice of thermometer. Because each (element, stage) density is a free unknown
(Section~\ref{sec:theory-saha}), the Saha fractions only \emph{scale} the design
columns, and a free per-stage density absorbs any scale. The search therefore sees
$T$ only through within-stage Boltzmann ratios---weak when the observed lines share
similar upper levels---and sees $\nel$ essentially not at all. Measured on a synthetic
Ca/Mg scene whose truth is \SI{10}{\kilo\kelvin}, the amplitude cost varies less than
\SI{1}{\percent} between $5.5$ and \SI{10}{\kilo\kelvin} and leans cold.

\alibz{} instead fixes $T$ by \emph{stage consistency}. At the true $(T,\nel)$ the
independent per-stage estimates of one element agree, because they measure the same
nuclei through the Saha fractions built into the columns; at a wrong state they do
not. For each element with at least two observed stages, \alibz{} takes the element's
own residual (observations minus every other species' fitted contribution) and fits
it twice---\emph{free}, one non-negative density per stage, and \emph{tied}, one
density through the summed stage columns---and charges the difference:
%
\begin{equation}
  \label{eq:tie-cost}
  C_\mathrm{tie}(T,\nel) =
    \sum_\mathrm{el}\max\!\bigl(0,\;
      \chi^2_\mathrm{tied} - \chi^2_\mathrm{free}\bigr),
\end{equation}
%
evaluated in the same squared units as the data misfit against the ungated design
matrix, with unit weight so the physics is trusted exactly as much as the data. On
the Ca/Mg synthetic the tie cost is $360$ at the truth node, against
$3.7\times10^3$ at the nearest point of the Saha degeneracy ridge and $2.6\times10^6$
at the state the amplitude objective alone preferred. Two-element thermometry leaves
a residual ridge along which $(T,\log_{10}\nel)$ trade off
($\SI{7}{\kilo\kelvin}/10^{15} \approx \SI{10}{\kilo\kelvin}/10^{17}
\approx \SI{17.5}{\kilo\kelvin}/10^{19}$); pairs with strong emissivity contrast
break it (the Si/K emissivity ratio changes $\sim$\num{2.5}$\times$ per
\SI{1000}{\kelvin}). A coarse $(T,\log_{10}\nel)$ grid seeds the optimizer by default,
and the reported composition remains the untied free solve, with $D_\mathrm{el}$ of
Equation~\eqref{eq:stage-disagreement} exposed as a diagnostic.

\subsection{The whole-pattern inverse problem}
\label{sec:theory-inverse}

The inverse problem is to recover concentrations and the plasma state from the full
pattern of fitted line areas, not from a handful of selected lines. Let
$\mathbf{y}$ be the vector of observed peak areas and let the design matrix
$\mathbf{A}(T,\nel)$ have one column per (element, stage) candidate, each column the
stacked emissivities of Equation~\eqref{eq:emissivity} (compressed by the escape
factor of Equation~\eqref{eq:escape} where self-absorption is active). At fixed
$(T,\nel)$ the concentrations $\mathbf{c}$ solve a weighted non-negative least-squares
(WNNLS) problem \citep{lawson1995nnls,bro1997nnls},
%
\begin{equation}
  \label{eq:wnnls}
  \hat{\mathbf{c}} = \argmin_{\mathbf{c}\ge 0}
    \sum_j \left(\frac{y_j - [\mathbf{A}\mathbf{c}]_j}{\sigma_{\mathrm{eff},j}}\right)^2,
\end{equation}
%
and an outer Gaussian-process optimizer \citep{jones1998ego} searches over
$(T,\nel,\sigma,\gamma)$ using the thermometer of Equation~\eqref{eq:tie-cost} (\texttt{peaky\_indexer\_v3}, built on
\citealp{harris2020numpy} and \citealp{virtanen2020scipy}). Non-negativity is the
physical statement that no line can have negative emission.

The effective per-peak variance is not the raw fit variance. \alibz{} floors it so
bright lines cannot claim fictitious precision:
%
\begin{equation}
  \label{eq:sigma-eff}
  \sigma_{\mathrm{eff},j}^2 =
    \sigma_{\mathrm{GLS},j}^2
    + (0.05\,y_j)^2
    + \sigma_\mathrm{floor}^2
    + \frac{y_j}{g},
\end{equation}
%
where $\sigma_{\mathrm{GLS},j}$ is the statistical sigma from the peak fit, the
\SI{5}{\percent} term is a fitting-procedure noise floor, $\sigma_\mathrm{floor}$ an
additive floor required because the generalized-least-squares sigma collapses on quiet
regions, and $y_j/g$ a provisional Poisson term (count gain $g=1$). The weighted solve
is opt-in; without supplied per-peak sigmas the model reduces to the historical
unweighted solve. On a synthetic Ca/Mg scene (truth $0.60/0.40$) the weighted solve at
the true state returns $0.611/0.389$, against $0.764/0.236$ unweighted. Because the
measured noise model is dominated by model inadequacy---the fit explains only
$\sim$\SI{9}{\percent} of $\chi^2$ on synthetics---absolute $r^2$ is not a reliable
score, and \alibz{} compares competing basins by $\Delta\chi^2$ and $\Delta$BIC
\citep{schwarz1978bic} rather than by $r^2$.

Per-element totals are formed by a tied non-negative one-dimensional re-solve of each
element's stages against its own residual; when the stages have disjoint support this
reduces exactly to the $\lVert\mathbf{A}_{:,s}\rVert^2$ precision-weighted average of
the per-stage estimates, which is well defined even when stages blend and avoids the
arbitrary vertex split that averaging raw NNLS coefficients would inherit.

Three further guards make the solve honest about what it can and cannot claim. First,
\emph{detectability}: an element enters the composition only if its strongest predicted
line clears
%
\begin{equation}
  \label{eq:detect}
  z = \max_j \frac{[\mathbf{A}_{:,\mathrm{el}}\,\hat{\mathbf{c}}]_j}{\sigma_{\mathrm{eff},j}}
    \ge 3,
\end{equation}
%
the score-test statistic for $c_\mathrm{el}>0$ against the local per-peak area noise
(\texttt{ELEMENT\_DETECT\_Z\_MIN}). A logistic presence probability and a regularized
(ridge or spike-and-slab \citep{mitchell1988spikeslab,george1993ssvs}) solve are proposed
but not built. Second, a
\emph{basin-collapse guard}: a degenerate single-element basin can score better on raw
$r^2$ than a plausible mixture (measured: a Bi-only vertex scored $r^2=0.53$ against a
plausible fit at $0.0$), so \alibz{} rejects a re-solve whose top element fraction
newly reaches \texttt{COLLAPSE\_TOP\_FRACTION}$=0.9$.

Third, a \emph{confounder contest} addresses near-collinear design columns, where the
NNLS vertex hands all the flux in a shared window to one of two elements and the
subsequent pruning makes the alternative unrepresentable---the error then certifies
itself by dragging the fitted temperature. For each element $E$ with support and each
rival $E'$, \alibz{} caps the rival by its own true negatives,
%
\begin{equation}
  \label{eq:contest}
  c_\mathrm{max}(E') =
    \min_j \frac{y_j + 3\sigma_j}{T_j(E')},
\end{equation}
%
over $E'$'s non-resonance lines only ($E_i>\SI{0.2}{\electronvolt}$, because resonance
lines self-absorb and would understate the rival), where $T_j(E')$ is the rival's
template response at peak $j$. A supporting peak of $E$ is \emph{contested} if some
rival at its cap still covers at least half of it. The scan runs existentially over a
corpus-plausible grid ($T\in\{6,8,10,12\}$\,kK $\times\;\log_{10}\nel\in\{16.5,17.5\}$)
against the pre-prune candidate table, never only at the fitted state. An element all
of whose support is contested is demoted to \texttt{confounded}; resolution credits it
only its clear flux and reattributes the contested remainder to the rival that can
globally host it, with the response ratio clamped at $1$ (no mass creation) and bounded
by the rival's own evidence. On a drill-core sample this moved Mn from
\SI{53}{\percent} to $0$ and Mg from \SI{8}{\percent} to \SI{31}{\percent}, while
genuinely independent detections were untouched. The full joint re-solve of a many-way
confounded group remains future work.
